\bookchapter{The debug server and robotcontrol}\label{k:rc}

You cannot fix what you cannot see. The robot has no screen, and \texttt{print} over SSH stops being useful the moment two loops print at once. This chapter builds a \textbf{debug server} in the robot code that streams the contents of the blackboard to a desktop program, \textbf{robotcontrol}, and receives commands from it. The messages are defined with \textbf{protobuf} in one \texttt{.proto} file, so that the team's other tools (for example tools written in Rust) can read the same stream. Milestone: one.

\begin{needbox}
\begin{itemize}[nosep]
  \item The two loops (Chapter~\ref{k:loops}) running against the fake robot. Python packages \texttt{protobuf} and \texttt{grpcio-tools} (only for generating code).
  \item Background: \rtr{Networking: From Cable to Socket}, \rtr{Protobuf Serialization}, \rtr{UDP Networking}.
\end{itemize}
\end{needbox}

\section*{The messages}
\begin{verbatim}
syntax = "proto3";
package k1debug;

message Vec3 { float x = 1; float y = 2; float z = 3; }
message BallPercept { bool seen = 1; float x = 2; float y = 3; }
message RobotPose { float x = 1; float y = 2; float theta = 3; }

message MotionCommand {
  string kind = 1;            // "stand", "walk" or "kick"
  float vx = 2; float vy = 3; float vyaw = 4;
  float ttl = 5;
  int32 priority = 6;
}

message Frame {               // robot -> robotcontrol
  uint32 cycle = 1;
  double time = 2;
  string game_state = 3;
  Vec3 gravity = 4;
  BallPercept ball = 5;
  RobotPose pose = 6;
  MotionCommand command = 7;  // what motion is executing
  map<string, string> extra = 15;  // anything else, as text
}

message Request {             // robotcontrol -> robot
  oneof what {
    MotionCommand command = 1;
    string subscribe = 2;     // e.g. "frames@10" for 10 frames per second
  }
}
\end{verbatim}
The important representations get their own typed messages; everything else can be sent as text in \texttt{extra} until it deserves a message. Field numbers are forever: never reuse a number for another meaning, or old logs and other tools will misread it. Generate the Python code with
\begin{verbatim}
python -m grpc_tools.protoc -Iproto --python_out=src/k1code proto/debug.proto
\end{verbatim}

\section*{Transport: TCP with a length prefix}
A TCP connection is a stream of bytes, not of messages, so each message is sent with its length in front (four bytes, little-endian):
\begin{verbatim}
import struct

def send_msg(sock, msg):
    data = msg.SerializeToString()
    sock.sendall(struct.pack("<I", len(data)) + data)

def recv_exact(sock, n):
    buf = b""
    while len(buf) < n:
        part = sock.recv(n - len(buf))
        if not part:
            raise ConnectionError("closed")
        buf += part
    return buf

def recv_msg(sock, cls):
    (n,) = struct.unpack("<I", recv_exact(sock, 4))
    msg = cls()
    msg.ParseFromString(recv_exact(sock, n))
    return msg
\end{verbatim}
Why TCP and not UDP? robotcontrol must not lose a command, and the stream for one viewer at $10$ frames per second is small: a frame with a ball, a game state and one extra value was $45$ bytes in my test.

\section*{The debug server}
\begin{verbatim}
class DebugServer:
    """Streams frames to one robotcontrol and passes its commands on."""

    def __init__(self, slot, port=5555):
        self.slot = slot                       # the CommandSlot of motion
        self.srv = socket.create_server(("0.0.0.0", port))
        self.conn = None
        threading.Thread(target=self._accept, daemon=True).start()

    def _accept(self):
        while True:
            conn, addr = self.srv.accept()
            self.conn = conn
            threading.Thread(target=self._read, args=(conn,),
                             daemon=True).start()

    def _read(self, conn):
        try:
            while True:
                req = recv_msg(conn, pb.Request)
                if req.WhichOneof("what") == "command":
                    c = req.command
                    self.slot.put(MotionCommand(
                        c.kind, c.vx, c.vy, c.vyaw, t=time.monotonic(),
                        ttl=c.ttl, source="robotcontrol",
                        priority=c.priority))
        except (ConnectionError, OSError):
            if self.conn is conn:
                self.conn = None

    def publish(self, frame):
        """Called by cognition once per cycle (or every n-th cycle)."""
        conn = self.conn
        if conn is None:
            return
        try:
            send_msg(conn, frame)
        except OSError:
            self.conn = None
\end{verbatim}
Commands from robotcontrol go straight into the \texttt{CommandSlot} of Chapter~\ref{k:loops}, with the source \texttt{robotcontrol} and the time of \emph{arrival} (the laptop's clock is not the robot's clock). I tested server and client on one laptop: a walk command with priority $10$ arrived in the slot and won; $0.6$\,s later it had expired and the slot returned \emph{stand}. The value $0.2$ came back as \texttt{0.20000000298023224}: protobuf's \texttt{float} has $32$ bits.

\section*{robotcontrol, version 1}
The first robotcontrol is a command-line program on the laptop:
\begin{verbatim}
uv run robotcontrol watch k1-03         # prints frames
uv run robotcontrol send k1-03 stand    # one command
\end{verbatim}
It looks up the robot's address in \texttt{config/robots/k1-03.toml}. Later versions add a live plot of selected values, a field view with pose and ball, and the keyboard driving of Chapter~\ref{k:cmd}. Because the stream is protobuf, a team member can write a faster or prettier viewer in another language from the same \texttt{.proto} file.

\section*{Steps}
\begin{stepbox}{Step 1: proto and code generation}
Add \texttt{proto/debug.proto}, a script that generates the Python module, and a test that builds a \texttt{Frame}, serializes and parses it.
\end{stepbox}
\begin{stepbox}{Step 2: a FrameBuilder module}
A cognition module that \texttt{requires} the representations you want to see and fills a \texttt{Frame}; it calls \verb!server.publish! every second cycle ($10$ frames per second).
\end{stepbox}
\begin{stepbox}{Step 3: robotcontrol watch and send}
Write both subcommands. Run the robot code with the fake robot, connect, watch the pose change while you send a walk command.
\end{stepbox}
\begin{stepbox}{Step 4: on the robot}
Deploy, connect from the laptop over the network, and repeat Step~3 with the robot hanging and the walk command \emph{not yet executed} (the output adapter only logs commands until Chapter~\ref{k:walk}).
\end{stepbox}

\begin{checkbox}
\begin{itemize}[nosep]
  \item robotcontrol shows $10$ frames per second from the fake robot and from the real one.
  \item Closing robotcontrol while it sends a walk command: the robot code falls back to cognition's command (or \emph{stand}) within the time to live.
\end{itemize}
\end{checkbox}

\begin{troublebox}
\begin{description}[leftmargin=0.6cm,nosep]
  \item[\texttt{Connection refused}] The server is not running, or a firewall blocks the port. \verb!ss -ltnp | grep 5555! on the robot.
  \item[Garbage values after a while] A message was read partly; use \verb!recv_exact!, never a single \verb!recv!.
  \item[Cognition becomes slow when the Wi-Fi is bad] \verb!sendall! blocks when the network is slow. Let \verb!publish! put the frame into a queue of length $1$ (dropping an unsent older frame) and send from a separate thread.
\end{description}
\end{troublebox}

\begin{keyidea}[What to remember]
\begin{itemize}[nosep]
  \item One \texttt{.proto} file defines the debug messages for all tools; field numbers never change meaning.
  \item TCP with a length prefix; commands go into the same command slot as cognition's, with source, priority and time to live.
  \item The debug stream must never slow down the robot.
\end{itemize}
\end{keyidea}
